\section{Conclusion}

NORA formulates long-tailed annular weld inspection as a query-level evidence-reliability problem rather than a score-thresholding problem. Its central design combines normal-reference IBO transitions in the spatial and local-frequency domains using candidate-specific reliability. The disagreement between domains is retained as evidence instead of being removed by a fixed feature mixture. On the industrial data, geometry preservation improved candidate availability and evidence-based re-evaluation reduced misses. The LVIS pilot showed that normal-reference reliability transfers at the candidate level. The conservative confusion corrector left aggregate LVIS instance statistics unchanged, which defines its present limitation. A completed B0--B6 end-to-end fusion ablation and an official LVIS AP evaluation are still required before claiming that the proposed fusion outperforms conventional spatial--frequency fusion. Subsequent validation should also assess reference-bank updates across acquisition conditions and the runtime cost of joint fusion in production inspection.
